\begin{onehalfspace}

L'implementazione segue un unico contratto: preset rumorosi controllati per $N=15$,
validazioni ideali separate per $N=21/35$, confronto appaiato delle strategie e
artefatti schema~2. I dettagli di netlist, manifest e organizzazione dei sorgenti sono
trasferiti nella Sezione~\ref{app:contratto_codice}.

% -----------------------------------------------
\section{L'algoritmo di Shor: implementazione di riferimento}
\label{sec:impl_shor}

Il circuito costruisce la \ac{QPE} con un registro di conteggio e uno di lavoro; applica
la moltiplicazione modulare controllata, la \ac{QFT} inversa e la misura. Il
post-processing ricava il periodo mediante frazioni continue e verifica i fattori non
banali con il \ac{GCD}. La QFT inversa è decomposta in Hadamard, rotazioni di fase
controllate e scambi; i relativi estratti sono nei
Listati~\ref{lst:circuito_completo_di_shor}
e~\ref{lst:implementazione_della_quantum_fourier}.

Per $N=15$, $a=7$, la rete textbook realizza realmente l'operazione elementare
$\times7\bmod15$ per il numero di potenze richiesto. Ciò corregge la precedente
selezione mediante il solo valore \texttt{power \% 4}, che non rappresentava in generale
$U^{\texttt{power}}$ (Listato~\ref{lst:modular_exponentiation_per_n15}). Per $N=21$ e
$N=35$ si impiega invece l'aritmetica reversibile di Beauregard, verificata idealmente
mediante truth table, pulizia degli ausiliari e prove QPE.

La compilazione, comune a training, baseline e sweep, fissa il livello di
ottimizzazione 2, il seed del transpiler e la base \texttt{rz/sx/x/cx}. Le
istanze maggiori sono escluse dalla campagna rumorosa completa per il costo della
netlist, ma non dai controlli ideali. Profondità, conteggi delle porte e hash SHA-256
delle compilazioni di riferimento sono riportati nella
Sezione~\ref{app:netlist_manifest}.

% -----------------------------------------------
\section{Configurazione del modello di rumore}
\label{sec:noise_model}

Il modello \texttt{qiskit\_aer.noise} è \textbf{uniforme e illustrativo}, non una
calibrazione hardware. Le \texttt{rz} sono cambi di frame virtuali, quindi senza durata
né errore; \texttt{sx/x} durano 50\,ns e \texttt{cx} 300\,ns. Depolarizzazione e
rilassamento termico agiscono su \texttt{sx/x} e \texttt{cx}, mentre il readout è
descritto da una matrice di confusione simmetrica separata
(Listato~\ref{lst:costruzione_del_noise_model}).

I simboli $\epsilon_{1q}$ e $\epsilon_{2q}$ indicano il parametro Aer $\lambda$ di
\texttt{depolarizing\_error}; per due qubit la probabilità Pauli non-identità è
$15\lambda/16$. Per i 224 gate \texttt{cx} della compilazione N15, il proxy che nessuno
di tali eventi si verifichi è ${\bigl(1-15\lambda/16\bigr)}^{224}$. Non è una probabilità di
successo: esclude rumore 1Q, rilassamento e readout.

% -----------------------------------------------
\section{Implementazione dei metodi proposti}
\label{sec:impl_metodi}

\subsection{Metodo 1: selezione TOP-1}
\label{sec:impl_metodo1}

M1 seleziona il primo esito per frequenza, tenta il post-processing e, in caso di
fallimento, passa all'iterazione successiva. Un tie-break SHA-256 derivato da seed e
bitstring rende l'ordinamento indipendente dal dizionario e dal valore numerico del
candidato (Listato~\ref{lst:pipeline_del_metodo_1}).

\subsection{Metodo 2: classificatore ML e ricerca TOP-4}
\label{sec:impl_metodo2}

M2 combina addestramento offline e inferenza online. La stessa netlist genera, da ogni
istogramma, le etichette TOP-1 e TOP-16; Random Forest, \ac{SVM} RBF e \ac{MLP} sono
confrontati con split stratificato. Il modello schema~2 registra definizione
dell'etichetta, revisione, preset, seed e hash del circuito, e l'inferenza rifiuta
manifest incompatibili. Se il filtro accetta l'istogramma, vengono provati fino a
quattro candidati: M1, TOP-4 e M2 ricevono pertanto gli stessi conteggi, tie-break e
post-processing. I Listati~\ref{lst:generazione_del_dataset_con}
e~\ref{lst:loop_di_inferenza_del} rendono verificabile il flusso; l'audit completo è
nella Sezione~\ref{app:classificatore}.

% -----------------------------------------------
\section{Infrastruttura di test e raccolta dati}
\label{sec:infrastruttura}

Le campagne N15 usano $K=30$ repliche appaiate e una sola simulazione per produrre
l'istogramma comune alle tre strategie. Seed sufficientemente distanziati evitano la
quasi sovrapposizione degli stream Aer; il primo successo è registrato separatamente e
un fallimento entro \texttt{max\_iter} riceve la sentinella
\texttt{max\_iter+1}. I confronti impiegano il Wilcoxon appaiato con
\texttt{zero\_method='pratt'}. Artefatti JSON, parametri, vettori appaiati, versioni,
hash e regole esatte per seed e test sono descritti nella
Sezione~\ref{app:contratto_codice} e sintetizzati dal
Listato~\ref{lst:schema_di_orchestrazione_degli}.

% -----------------------------------------------
\section{Implementazione della correzione d'errore}
\label{sec:impl_qec}

La seconda fase sostituisce l'istogramma finale di un circuito profondo con la sindrome
di circuiti di memoria. Qiskit Aer gestisce il codice a ripetizione, Steane e il proxy
di Shor; Stim e PyMatching rendono praticabile il surface code, mentre scikit-learn
serve il decoder appreso. La mappa completa fra moduli, strumenti e costo dominante è
nella Tabella~\ref{tab:organizzazione_qec} dell'Appendice~\ref{app:strumentiecodice}.

Il flusso comune prepara lo stato logico, inserisce il canale nel punto dichiarato,
estrae la sindrome su ancille e consegna gli eventi al decoder; l'esito è un successo o
fallimento logico, non una misura diretta della fedeltà di Shor. Repetition, Steane e
surface code realizzano questo schema con livelli crescenti di protezione. La curva di
sensibilità di Shor costituisce invece un secondo livello: riceve un proxy per gate e
misura la resa in fattori. Tenere separati i due livelli evita di convertire senza
giustificazione il tasso per memoria dei codici nel parametro per gate del circuito.

Nel codice a ripetizione la parità è copiata su ancille senza misurare i dati; nel codice
di Steane $\ket{0_L}$ è preparato e decodificato con una lookup table a peso minimo,
dichiarando ideale l'estrazione di sindrome. Per il surface code il crosstalk viene
iniettato anche dentro i blocchi \texttt{REPEAT} e \ac{MWPM} opera sui detection event.
Nel proxy di Shor, il $p_L$ per gate viene convertito nella convenzione depolarizzante di
Aer; non è identificato con il logical error rate per round senza un mapping aggiuntivo.
Il decoder appreso e \ac{MWPM} ricevono gli stessi eventi. Algoritmi, assunzioni e
listati per ciascun componente sono raccolti nella Sezione~\ref{app:dettagli_qec}.

Ogni script propaga il seed e registra nel JSON campioni, parametri e versioni; le figure
sono generate dagli artefatti senza trascrizione manuale. Questa disciplina costituisce
il contratto di riproducibilità comune alle due fasi.

\end{onehalfspace}
